\begin{pastexam}[label={px:gw3-115-ds-02}]{關務三等資訊處理 115 年 第二題}
（申論；20 分）

何謂 Complete Binary Tree？假設單一節點的樹其樹高為 1，證明若 Complete Binary Tree 含 n 個節點且樹高為 h，則
$2^{h-1} \le$ n $\le 2^h - 1$。（20 分）
\end{pastexam}

\begin{pastsoln}
\textbf{定義}（資結書第~7.1~節）。 樹高為 $h$ 的二元樹稱為完整二元樹（complete binary tree），如果：
\begin{enumerate}[label=(\arabic*)]
\item 第 $1$ 層到第 $h - 1$ 層都是滿的，第 $d$ 層恰有 $2^{d-1}$ 個節點；
\item 第 $h$ 層（最底層）的節點全部靠左排，中間沒有空缺。
\end{enumerate}
等價的說法：把高 $h$ 的滿二元樹由上到下、同一層由左到右編號 $1, 2, \ldots, 2^h - 1$，$n$ 個節點的完整二元樹就是
編號 $1$ 到 $n$ 的那些節點。 （有些課本把「每一層都滿」的樹也叫 complete；本題用的是上面這個常見的定義，滿二元樹是
它的特例。）

\textbf{引理}：滿的第 $d$ 層有 $2^{d-1}$ 個節點，第 $1$ 到第 $d$ 層全滿時共有 $2^d - 1$ 個節點。

\textbf{證明。} 第 $1$ 層只有根，$2^0 = 1$ 個。 第 $d$ 層滿時，每個節點恰有兩個子節點，第 $d + 1$ 層若滿就有
$2 \cdot 2^{d-1} = 2^d$ 個，由歸納法得第 $d$ 層 $2^{d-1}$ 個。 前 $d$ 層合計是等比級數
$1 + 2 + \cdots + 2^{d-1} = 2^d - 1$。 \hfill$\blacksquare$

\textbf{主要結論的證明。} 設完整二元樹有 $n$ 個節點、樹高 $h$（單一節點高 $1$，所以共有 $h$ 層）。 記第 $h$ 層的節點數
為 $k$。
\begin{itemize}
\item 由定義 (1)，第 $1$ 到第 $h - 1$ 層全滿，由引理共 $2^{h-1} - 1$ 個節點；所以 $n = (2^{h-1} - 1) + k$。
\item 下界：樹高是 $h$，第 $h$ 層至少有一個節點，$k \ge 1$，所以 $n \ge 2^{h-1} - 1 + 1 = 2^{h-1}$。
\item 上界：第 $h$ 層最多是滿的，$k \le 2^{h-1}$，所以 $n \le 2^{h-1} - 1 + 2^{h-1} = 2^h - 1$。
\end{itemize}
因此 $2^{h-1} \le n \le 2^h - 1$。 \hfill$\blacksquare$

兩個等號都取得到：只有最左邊一個節點在第 $h$ 層時 $n = 2^{h-1}$（例如 $h = 3$、$n = 4$）；第 $h$ 層滿時是滿二元樹，
$n = 2^h - 1$（$h = 3$、$n = 7$）。 把不等式取對數，$h - 1 \le \lg n < h$，得 $h = \lfloor \lg n \rfloor + 1$，這就是
資結書定理~7.3：$n$ 個節點的完整二元樹高度是 $\lfloor \lg n \rfloor + 1$，也是堆積的操作只要 $O(\log n)$ 的原因。

注意上界 $n \le 2^h - 1$ 對\textbf{任何}高 $h$ 的二元樹都成立（每層最多 $2^{d-1}$ 個），下界 $n \ge 2^{h-1}$ 才用到
「完整」：一般的二元樹高 $h$ 可能只有 $h$ 個節點（一條鏈）。
\end{pastsoln}
